\documentclass[11pt,a4paper]{article}
\usepackage[T1]{fontenc}
\usepackage[margin=22mm]{geometry}
\usepackage{isabelle,isabellesym}
\usepackage{amsmath,amssymb}
\usepackage{microtype}
\usepackage{pdfsetup}
\urlstyle{rm}
\isabellestyle{it}

\begin{document}
\emergencystretch=3em
\title{Cross-Chain Message Integrity}
\author{Jinwook Kim (Jay)}
\date{September 2026}
\maketitle

\begin{abstract}
This Isabelle/HOL development connects authenticated source facts to
destination-credit execution. Under explicit cryptographic, honest-attestation
and stable-source assumptions, accepted messages bind the full transfer
payload and arbitrary delivery traces create at most one new destination
credit per source event. A separate kernel characterization proves that
destination-local once markers give global uniqueness exactly when admitted
requests for the same source key have one destination. Constructive normal
and bypass decoders preserve rejection, effects and finite continuations.
A checked internal message summary characterizes the information distinguished
by the receiver over its declared contexts. Concrete regulatory consumers and
semantic counterexamples distinguish source authentication, current authority,
state preservation and recovery information. These are model-level results,
not proofs of cryptographic implementations, external finality, distributed
atomicity or deployed software.
\end{abstract}

\clearpage
\tableofcontents
\clearpage

\section{Scope and Use}

A message identifier names an envelope; a source key names the event whose
effect may be delivered again through another envelope or route. The binding
records the source key, asset, amount, intended destination, recipient,
historical source epoch, operation and domain separator. A destination credit
and its persistent once marker are added in one local model step.

The model counts execution occurrences. Two executions with identical values
count twice. Source debits, replicated records, repeated replies and later
lawful transfers are different operations and are not counted as new
destination credits. The global history is a proof observation; the receiver
checks the marker for its declared destination and source key.

\subsection{Assumptions and implementation obligations}

\begin{itemize}
\item Cryptographic verification is assumed to identify the actual signers of
the complete source statement. No BLS, hash or serialization implementation is
verified here.
\item The trusted epoch configuration supplies a finite roster, a Byzantine
fault bound and a threshold strictly above that bound. Honest roster members
attest true source facts. Stable source outcomes are uniquely associated with
their source keys. These are external source assumptions.
\item The execution context supplies the actual current authority, policy
permissions and state version. The caller input must be obtained from an
authenticated execution context, not accepted as a self-asserted wire identity.
Establishing these inputs for an implementation remains an obligation.
\item The local credit effect and its marker are recorded atomically and
durably. Reconstructing markers from a complete valid credit history is proved;
storage correctness and recovery of unrelated financial effects are not.
\item All modeled entry routes use the receiver. Exhaustive correspondence to
real contract, administrative, adapter and runtime entrypoints is not established.
\end{itemize}

There is no network liveness theorem, source-lock protocol, independent
terminal-decision protocol, full financial-journal recovery or implementation
refinement in this session. A source reversal certificate is distinguishable
from source finality authorizing a destination credit.

\section{Checked Results}

\paragraph{Authentication and global uniqueness.}
Threshold authentication obtains an honest signer from the roster and then
uses the separately stated source assumptions. The execution proof uses that
authenticated binding to connect destination-local markers with a global
source-event bound. A general kernel theorem supplies both directions of the
destination-functionality criterion; a two-message execution witnesses failure
when the condition is removed.

\paragraph{Representations and information.}
The normal-to-bypass encoding checks raw payload consistency before discarding
the redundant payload copy. Its decoder preserves complete replies and effects,
including rejection. The checked summary is computed after intrinsic
verification. Equality of summaries is equivalent to equality of the current
receiver's results over every context and state in the model. This is a semantic
information quotient, not an optimal byte encoding, a replacement wire
certificate or a characterization of every product-reachable policy.

\paragraph{Regulatory consumers.}
The transfer-permission consumer uses Regulatory Action Composition to keep
ordinary transfers and authorized enforcement distinct. A separate completed
state-update consumer uses Cross-Domain State Preservation for atomic
regulatory metadata updates. It preserves a source-key application marker.
Neither consumer identifies token supply with regulatory holding support.
The imported models remain within their existing scope~\cite{cdsp}.

\paragraph{Counterexamples and boundary witnesses.}
The session includes field-removal witnesses for the authenticated binding,
equal-value and dual-destination replay, marker loss across authority rotation,
untrusted threshold and repeated signer identifiers, raw-input laundering,
stale authority and version checks, source-reversal separation, and current
permission recovery. In particular, a stale FREEZE can preserve the CDSP state
invariant while violating the required authorization check.

\section{Relationship to Existing Results}

Durable completion records and consistent retry rendezvous are established
principles~\cite{rifl}. Atomic commit protocols separately specify stable,
non-conflicting decisions~\cite{graylamport}; Byzantine quorum intersection is
also established theory~\cite{malkhireiter}. The general factorization criterion
used here is standard. This development supplies checked statements and
counterexamples for the displayed receiver and source-fact profile. It does
not claim a new general atomic-commit protocol or a general theory of
distributed recovery.

\section{Reproduction}

Use Isabelle2025-2 with the pinned ADS\_Functor dependency recorded by the
repository. Build the named Cross\_Chain\_Message\_Integrity session with the
repository registered as a session directory. The repository README and
session README give the complete commands. The following theories are the
authoritative definitions, proofs and concrete witnesses.

\clearpage
\input{session}
\clearpage
\bibliographystyle{abbrv}
\bibliography{root}
\end{document}
